\subsection{FAMILIES OF FINITE NUMBERS OF ARBITRARY SIGN SUMMABLE IN R}

THEOREM 3. Let \((x_{t})_{t\in I}\) be a family of finite real numbers; then the following statements are equivalent:

\begin{enumerate}
\def\labelenumi{\alph{enumi})}
\item
  The family \((x_{i})\) is summable in \(\mathbf{R}\) .
\item
  The family \((|x_{i}|)\) is summable in \(\mathbf{R}\) .
\item
  The set of finite partial sums of the family \((x_{t})\) is bounded in R.
\end{enumerate}

Let \(I_1\) be the set of all \(i \in I\) such that \(x_i \geq 0\) , and \(I_2\) the set of all \(i \in I\) such that \(x_i < 0\) . The family \((x_i)_{i \in I}\) {[}resp. \((|x_i|)_{i \in I}\) {]} is summable if and only if each of the families \((x_i)_{i \in I_1}\) and \((x_i)_{i \in I_2}\) {[}resp. \((|x_i|)_{i \in I_3}\) and \((|x_i|)_{i \in I_4}\) {]} is summable (Chapter III, § 5, no. 3, Propositions 2 and 3). Now it comes to the same thing to say that \((x_i)_{i \in I_1}\) is summable, or that \((|x_i|)_{i \in I_2}\) is summable, or that the set of finite partial sums of the family \((x_{i})_{i\in I_1}\) is bounded (no. 1, Theorem 1); and the same is true with \(I_1\) replaced by \(I_2\) . The theorem follows immediately.

Theorem 3 shows that the summability in R of a family of finite real numbers depends only on the summability of the family of their absolute values.

We recall (Chapter III, § 5, no. 5, Proposition 6) that if \((x_{t})\) and \((y_{t})\) are two summable families of finite real numbers, then the family \((x_{t} + y_{t})\) is summable and \(\sum_{i}(x_{t} + y_{t}) = \sum_{i}x_{t} + \sum_{i}y_{t}\) . Moreover, if \((x_{t})\) is a summable family of finite real numbers and \(a\) is any finite real number, then the family \((ax_{t})\) is summable in \(\mathbf{R}\) , and we have \(\sum_{i}ax_{t} = a\) . \(\sum_{i}x_{t}\) .

